\documentclass[11pt,a4paper]{article}

% =========================================================
%  STAR TOKAMAK MONOGRAPH SECTION -- DETAILED ENGINEERING VERSION
% =========================================================

% XeLaTeX
\usepackage{fontspec}
\setmainfont{TeX Gyre Termes}

% Idioma
\usepackage[spanish,es-nodecimaldot]{babel}

% Márgenes
\usepackage[margin=2.5cm]{geometry}

% Matemáticas y símbolos
\usepackage{amsmath,amssymb,bm}

% Figuras y tablas
\usepackage{graphicx}
\usepackage{booktabs}
\usepackage{tabularx}
\usepackage{array}
\usepackage{longtable}
\usepackage{multirow}
\usepackage{subcaption}
\usepackage{float}
\usepackage{xcolor}

% Hipervínculos
\usepackage[colorlinks=true,linkcolor=blue,citecolor=blue,urlcolor=blue]{hyperref}

% Mejor espaciado en tablas
\renewcommand{\arraystretch}{1.15}

\begin{document}

% =========================================================
% SECCIÓN: STAR TOKAMAK
% =========================================================

\section{STAR Tokamak}
\label{sec:star_tokamak}


\subsection{Introducción al concepto STAR}
El \emph{Spherical Tokamak Advanced Reactor} (STAR) es un concepto preconceptual de planta
fusionada tipo tokamak esférico desarrollado con una filosofía distinta a la de muchos diseños
puramente centrados en desempeño de plasma. En lugar de partir únicamente de una configuración
magnética deseada y posteriormente adaptar la ingeniería alrededor de ella, STAR se plantea como
una arquitectura integrada en la que la física del plasma, la disposición de bobinas, la segmentación
del blanket, el acceso para mantenimiento, el costo de construcción y la disponibilidad operacional
se consideran simultáneamente. Esta motivación es central para el presente trabajo, porque el
equilibrio magnético no puede evaluarse de forma aislada: debe coexistir con una geometría de
primera pared, blanket, estructuras pasivas y bobinas que respete restricciones reales de integración.

Desde el punto de vista geométrico, STAR pertenece a la familia de tokamaks esféricos. Esto implica
un aspect ratio bajo, típicamente alrededor de $A=R_0/a\approx 2$, y una sección poloidal más
compacta, elongada y de alto shaping que la de un tokamak convencional de mayor aspect ratio. En
el diseño de referencia, esta compacidad se explota como una posible vía para reducir tamaño de
reactor, costo de capital y complejidad de componentes, sin renunciar a una configuración de plasma
de alto desempeño. En consecuencia, la forma del plasma buscada en este proyecto no fue una
sección circular ni una elipse simple, sino una configuración tipo D, con elongación elevada,
triangularidad positiva y topología diverted.

La arquitectura STAR también introduce restricciones de ingeniería particularmente importantes.
El diseño de referencia considera una estrategia de mantenimiento vertical, una reducción del número
de bobinas de campo toroidal a 12, sectores de blanket DCLL colapsables y una disposición pensada
para mejorar disponibilidad y simplificar el reemplazo de componentes internos. En el contexto de
esta monografía, estos elementos justifican la decisión de no tratar el dominio entre plasma y bobinas
como un espacio vacío arbitrario. Por el contrario, la sección poloidal debe distinguir al menos cuatro
entidades: bobinas activas PF/CS, primera pared o limiter, región de blanket y estructuras pasivas
asociadas al vessel/soporte. Esta separación fue precisamente uno de los cambios que permitió que el
modelo numérico adquiriera consistencia ingenieril.

El objetivo de esta sección, por tanto, no es reproducir exhaustivamente el diseño reactor completo,
sino construir una versión axisimétrica y computacionalmente manejable del problema de equilibrio
que conserve las restricciones geométricas más relevantes. El modelo desarrollado parte de un CAD
poloidal, identifica familias de bobinas, asigna estructuras pasivas, impone una frontera de plasma
accesible y resuelve el equilibrio Grad--Shafranov free-boundary con FreeGSNKE. Bajo esta lectura,
STAR no se usa solamente como inspiración geométrica, sino como un caso de estudio de diseño
ingenieril: el equilibrio aceptable debe ser compatible con la forma del plasma, la arquitectura del
blanket, la ubicación del central solenoid, las bobinas PF y los márgenes de corriente.

\subsection{Propósito de esta sección}
Esta sección presenta el desarrollo de un modelo de equilibrio free-boundary Grad--Shafranov
para un tokamak esférico tipo STAR, con énfasis en el proceso de diseño ingenieril que llevó
al equilibrio baseline actual. El objetivo no es únicamente mostrar una solución final, sino
documentar de forma reproducible cómo se construyó la máquina numérica a partir de CAD,
cómo se interpretaron las estructuras físicas de la sección poloidal, cómo se depuró la
discretización de bobinas y cómo se integró el análisis de separatriz, X-points, contención,
corrientes y perfiles magnéticos.

A diferencia de iteraciones anteriores, el modelo actual ya no utiliza el plasma target del CAD
como restricción dominante. La geometría del plasma finalmente obtenida se considera más
adecuada para la configuración electromagnética actual que la curva objetivo inicial. Por ello,
el presente reporte trata al antiguo plasma target como una referencia histórica útil para
exploración, pero no como la definición final de éxito. El nuevo criterio de aceptación combina:

\begin{itemize}
    \item forma D físicamente plausible;
    \item topología double-null;
    \item LCFS confinada dentro de la primera pared;
    \item legs de separatriz dirigidas hacia las regiones de divertor;
    \item corrientes de bobinas dentro de límites ingenieriles razonables;
    \item diagnósticos magnéticos y geométricos consistentes.
\end{itemize}

El equilibrio actual debe entenderse como un \emph{baseline geométrico de baja presión}. Ya
proporciona una forma STAR-like defendible, pero todavía no representa el punto final del
reactor a alta presión. La siguiente etapa será continuar esta rama en presión y recalcular
los diagnósticos MHD asociados.

\subsection{Resumen ejecutivo del resultado alcanzado}
El resultado principal de esta etapa es un equilibrio double-null confinado, con forma D,
aspect ratio cercano a 2, elongación mayor a 2 y triangularidad positiva cercana a 0.47. Los
parámetros principales del baseline son:

\begin{table}[htbp]
\centering
\caption{Resumen del equilibrio baseline actual.}
\label{tab:baseline_summary_exec}
\begin{tabular}{lc}
\toprule
\textbf{Magnitud} & \textbf{Valor} \\
\midrule
Corriente de plasma, $I_p$ & $4.0~\mathrm{MA}$ \\
Presión central usada en el baseline, $p_{\mathrm{axis}}$ & $2.0\times 10^3~\mathrm{Pa}$ \\
$R_{\mathrm{axis}}$ & $4.397~\mathrm{m}$ \\
$Z_{\mathrm{axis}}$ & $\approx 0~\mathrm{m}$ \\
$R_0$ & $4.208~\mathrm{m}$ \\
$a$ & $2.076~\mathrm{m}$ \\
$A=R_0/a$ & $2.028$ \\
$\kappa$ & $2.163$ \\
$\delta_u,\delta_l$ & $0.469,0.469$ \\
Área poloidal & $27.47~\mathrm{m^2}$ \\
Volumen toroidal aproximado & $726.27~\mathrm{m^3}$ \\
Distancia mínima LCFS--\texttt{WALL\_INNER} & $9.78~\mathrm{cm}$ \\
Topología & Double-null reconstruida \\
\bottomrule
\end{tabular}
\end{table}

El resultado es importante porque se obtuvo sin introducir bobinas de shaping ficticias fuera
del coilset del CAD. La mejora decisiva no provino de aumentar artificialmente la complejidad
del sistema de control, sino de corregir la interpretación geométrica de la máquina: el central
solenoid fue reposicionado de acuerdo con el gap observado en el diseño de referencia, las
paredes internas y externas se reinterpretaron como first wall y blanket back plate, y la capa
\texttt{STAR\_VESSEL} se modeló como estructura pasiva.

\subsection{Evolución del modelo y lecciones de las iteraciones fallidas}
La primera etapa del proyecto consistió en construir una máquina STAR-like simplificada,
resolver equilibrios axisimétricos y obtener configuraciones visualmente plausibles. Sin embargo,
conforme se incorporó la geometría CAD realista, se hizo evidente que el problema era más
restrictivo que un ajuste de parámetros geométricos. Tres dificultades dominaron las primeras
iteraciones:

\begin{enumerate}
    \item \textbf{El plasma target original no era necesariamente alcanzable.} La curva objetivo
    inicial era útil como guía, pero no estaba garantizado que correspondiera a una rama de
    equilibrio compatible con la disposición realista de bobinas, paredes y divertores.

    \item \textbf{La semántica de paredes y estructuras estaba mezclada.} En iteraciones tempranas,
    \texttt{WALL\_INNER}, \texttt{WALL\_OUTER}, blanket y vessel se trataban de forma demasiado
    simplificada. Esto permitía aceptar soluciones que en realidad no respetaban la región
    accesible al plasma.

    \item \textbf{La autoridad del central solenoid estaba mal representada.} El CS se encontraba
    demasiado cercano a $R=0$ en el modelo inicial. Aun con corrientes elevadas, su contribución
    al shaping no se comportaba como se esperaba. Después de corregir el gap radial, el CS empezó
    a influir en la forma del plasma de manera mucho más física.
\end{enumerate}

Estas iteraciones no se consideran resultados finales, pero fueron necesarias para identificar
el origen del problema. El desarrollo siguió una lógica de depuración ingenieril: cuando una
optimización fallaba o producía formas espurias, se revisaba si el fallo provenía del solver,
de la función objetivo, de la interpretación CAD, de la discretización de bobinas o de la
clasificación de separatriz. Esta estrategia permitió separar problemas numéricos de problemas
físicos de modelado.

\begin{figure}[htbp]
\centering
\fbox{%
\parbox{0.92\linewidth}{%
\centering
\textbf{[INSERTAR FIGURA]}\\[2mm]
Comparación conceptual entre: (a) una geometría inicial simplificada o mal interpretada, (b)
un equilibrio con inboard bulge o contención deficiente, y (c) el baseline actual después de
corregir semántica CAD y posición del CS.%
}}
\caption{Evolución cualitativa del modelo desde las iteraciones exploratorias hasta el baseline actual.}
\label{fig:star_model_evolution_placeholder}
\end{figure}

\subsection{Arquitectura computacional del pipeline}
El pipeline final se organizó como una cadena modular:

\[
\mathrm{CAD}\rightarrow
\mathrm{machine\ builder}\rightarrow
\mathrm{equilibrium\ solver}\rightarrow
\mathrm{topology\ analysis}\rightarrow
\mathrm{diagnostics}\rightarrow
\mathrm{documentation}.
\]

Cada módulo cumple una función específica, de modo que el modelo pueda depurarse por etapas.
La separación de responsabilidades fue esencial para no confundir errores de lectura CAD con
problemas de equilibrio o con fallos del post-proceso.

\subsubsection{\texttt{star\_machine\_cad.py}: constructor de máquina desde CAD}
El módulo \texttt{star\_machine\_cad.py} importa el DXF y produce una representación numérica de
la máquina. Sus tareas principales son:

\begin{itemize}
    \item leer polilíneas y polígonos por capas CAD;
    \item reconstruir paredes internas y externas;
    \item extraer bobinas activas \texttt{COIL\_*};
    \item discretizar bobinas en filamentos distribuidos espacialmente;
    \item agrupar bobinas por familia física: \texttt{CS}, \texttt{CS\_MID}, \texttt{CS\_END},
    \texttt{PF1}, \ldots, \texttt{PF6};
    \item calcular pesos de distribución de corriente por área;
    \item calcular límites de corriente recomendados por familia;
    \item importar \texttt{STAR\_VESSEL} como estructuras pasivas;
    \item guardar metadata de materiales, resistividad y semántica geométrica.
\end{itemize}

El cambio más importante respecto al modelo anterior fue desactivar el relleno artificial del
blanket. Antes se generaban filamentos pasivos automáticos entre \texttt{WALL\_INNER} y
\texttt{WALL\_OUTER}; esto era útil como aproximación temprana, pero podía introducir metal
ficticio en regiones donde la geometría real del CAD ya estaba disponible. En el modelo actual,
la estructura pasiva principal proviene de \texttt{STAR\_VESSEL}.

\subsubsection{\texttt{star\_equilibrium.py}: equilibrio free-boundary}
El módulo \texttt{star\_equilibrium.py} construye el objeto de máquina, aplica corrientes por
familias y resuelve la ecuación de Grad--Shafranov con FreeGSNKE. La formulación usada en el
baseline emplea un perfil tipo \texttt{ConstrainPaxisIp}, parametrizado por:

\[
\left\{p_{\mathrm{axis}}, I_p, f_{\mathrm{vac}}, \alpha_m, \alpha_n\right\}.
\]

Para el baseline actual se usó $I_p=4.0~\mathrm{MA}$, $p_{\mathrm{axis}}=2.0\times10^3~\mathrm{Pa}$
y $f_{\mathrm{vac}}=20.8~\mathrm{T\,m}$. Estos valores no representan todavía el estado final de
presión del reactor; se eligieron para construir una rama estable de baja presión y validar la
geometría.

La solución se obtiene mediante una rampa de continuation interna. En lugar de aplicar todas
las corrientes de golpe, el solver escala progresivamente las corrientes con factores $f$ hasta
llegar al estado final. Esta estrategia reduce fallos de convergencia y evita que el solver salte
a ramas espurias durante las primeras iteraciones.

\subsubsection{\texttt{analyze\_star.py}: reconstrucción de LCFS y topología}
El módulo \texttt{analyze\_star.py} extrae la geometría del plasma y clasifica la topología.
Durante el proyecto se descubrió que contornear exactamente en $\psi_x$ podía fallar cerca de
los X-points, porque la separatriz es una curva singular. Por esta razón se implementó una
estrategia de reconstrucción más robusta:

\begin{enumerate}
    \item localizar el eje magnético y los puntos críticos candidatos;
    \item detectar X-points superiores e inferiores;
    \item evaluar los segmentos de contorno asociados al nivel de separatriz;
    \item reconstruir una curva LCFS cerrada o near-separatrix válida para métricas geométricas;
    \item calcular $R_0$, $a$, $A$, $\kappa$, $\delta_u$, $\delta_l$, área, perímetro y bounds;
    \item verificar contención respecto a \texttt{WALL\_INNER}.
\end{enumerate}

Esta reconstrucción fue particularmente importante para configuraciones double-null, donde
los segmentos de separatriz pueden aparecer como ramas abiertas en el contorno exacto, pero
la frontera cerrada del plasma se reconstruye de forma consistente a partir de los segmentos
físicamente conectados.

\subsubsection{\texttt{diagnose\_star\_equilibrium.py}: diagnósticos tipo paper}
El script \texttt{diagnose\_star\_equilibrium.py} genera tablas y figuras para documentación.
Produce, entre otros archivos, un resumen en JSON, tablas CSV y figuras de equilibrio,
perfiles de campo magnético, perfil de presión, safety factor, mapa de corriente toroidal y
utilización de bobinas.

El propósito de este script es separar el cálculo del equilibrio de su análisis documental. De
este modo, una misma solución puede diagnosticarse de forma reproducible y compararse contra
otras ramas futuras cuando se continúe en presión.

\begin{table}[htbp]
\centering
\caption{Resumen de módulos principales del pipeline.}
\label{tab:pipeline_modules_detailed}
\begin{tabularx}{\textwidth}{p{0.28\textwidth}X}
\toprule
\textbf{Módulo} & \textbf{Función dentro del pipeline} \\
\midrule
\texttt{star\_machine\_cad.py} & Importa el DXF, interpreta capas, discretiza bobinas activas y estructuras pasivas, calcula grupos y límites de corriente. \\
\texttt{star\_equilibrium.py} & Construye el equilibrio free-boundary, aplica corrientes de bobinas, ejecuta continuation y produce plots base. \\
\texttt{analyze\_star.py} & Reconstruye LCFS/separatriz, X-points, métricas geométricas y contención contra la first wall. \\
\texttt{diagnose\_star\_equilibrium.py} & Genera diagnósticos tipo paper: geometría, perfiles magnéticos, presión, $q$, corrientes y tablas. \\
\texttt{fit\_simplified\_dn\_*.py} & Scripts de optimización usados durante exploración. Actualmente se conservan como herramientas auxiliares, no como definición final del baseline. \\
\bottomrule
\end{tabularx}
\end{table}

\subsection{Interpretación física de las capas CAD}
La interpretación final del CAD quedó definida de la siguiente manera:

\begin{center}
\begin{tabular}{ll}
\toprule
\textbf{Capa CAD} & \textbf{Interpretación física/numerical} \\
\midrule
\texttt{COIL\_*} & Bobinas activas PF/CS \\
\texttt{STAR\_VESSEL} & Estructura pasiva metálica de soporte/vessel \\
\texttt{WALL\_INNER} & First wall / limiter / frontera accesible al plasma \\
\texttt{WALL\_OUTER} & Frontera externa del blanket / back plate \\
\texttt{PLASMA\_TARGET} & Referencia histórica, no objetivo final principal \\
\bottomrule
\end{tabular}
\end{center}

La jerarquía geométrica resultante puede describirse como:

\[
\mathrm{PF/CS\ coils}\;\supset\;\mathrm{STAR\_VESSEL}\;\supset\;\mathrm{WALL\_OUTER}\;\supset\;\mathrm{WALL\_INNER}\;\supset\;\mathrm{plasma\ region}.
\]

En esta convención, \texttt{WALL\_INNER} es la frontera más importante para el análisis de
plasma: si la LCFS sale de \texttt{WALL\_INNER}, el plasma invade la primera pared o el blanket,
por lo que la solución debe considerarse físicamente inválida. \texttt{WALL\_OUTER} no define
la región accesible al plasma; representa la frontera externa del blanket. \texttt{STAR\_VESSEL}
se usa como estructura pasiva externa al blanket.

\begin{figure}[htbp]
\centering
\fbox{%
\parbox{0.92\linewidth}{%
\centering
\textbf{[INSERTAR FIGURA]}\\[2mm]
Plot limpio del setup de máquina: bobinas activas por familia, \texttt{STAR\_VESSEL} en gris,
\texttt{WALL\_OUTER} como frontera externa del blanket y \texttt{WALL\_INNER} como limiter.%
}}
\caption{Semántica final de las capas CAD usadas en el modelo numérico.}
\label{fig:cad_semantics_final_placeholder}
\end{figure}

\subsection{Discretización de bobinas activas}
Una fuente inicial de error fue que las bobinas del CAD se interpretaban de forma demasiado
concentrada, parecida a corrientes puntuales o a un número insuficiente de centros de filamento.
Esto era problemático porque el campo de una bobina extensa no debe representarse como una
fuente puntual en su centro geométrico cuando el plasma está cerca de la bobina.

El método corregido discretiza cada bobina rectangular o poligonal en una malla de filamentos.
Cada filamento representa una fracción del área conductora y recibe una fracción proporcional
de la corriente total de su familia. Para una familia $F$ con filamentos $i=1,\ldots,N_F$, se
definió el peso:

\[
w_i = \frac{A_i}{\sum_{j=1}^{N_F} A_j},
\]

de modo que la corriente asignada a cada filamento sea:

\[
I_i = w_i I_F.
\]

Este procedimiento garantiza que:

\[
\sum_{i=1}^{N_F} I_i = I_F.
\]

El diagnóstico de discretización confirmó que no existían centros duplicados y que la distancia
mínima entre centros era finita. En el modelo final se obtuvieron $602$ filamentos activos.
Las familias activas fueron:

\begin{table}[htbp]
\centering
\caption{Número de filamentos activos por familia en el modelo final.}
\label{tab:active_filaments_by_family}
\begin{tabular}{lc}
\toprule
\textbf{Familia} & \textbf{Número de filamentos} \\
\midrule
PF1 & 90 \\
PF2 & 90 \\
PF3 & 50 \\
PF4 & 36 \\
PF5 & 48 \\
PF6 & 96 \\
CS & 192 \\
CS\_MID & 88 \\
CS\_END & 104 \\
\bottomrule
\end{tabular}
\end{table}

La segmentación \texttt{CS\_MID}/\texttt{CS\_END} no introduce bobinas externas nuevas; representa
una descomposición del winding pack central para estudiar su contribución al shaping. En la
interpretación final, la corriente total del CS debe analizarse tanto de forma global como por
regiones locales del winding pack.

\subsection{Estructuras pasivas y materiales}
La capa \texttt{STAR\_VESSEL} se discretizó como estructura pasiva metálica. A diferencia de las
bobinas activas, estos filamentos se inicializan con corriente cero y no se incluyen como
familias de control. Su función principal en el modelo estático es documentar la geometría
física de la máquina y preparar el pipeline para extensiones dinámicas o resistivas.

El modelo usa las siguientes asignaciones de material como metadata ingenieril:

\begin{table}[htbp]
\centering
\caption{Materiales y resistividades efectivas usadas como metadata del modelo.}
\label{tab:materials_metadata}
\begin{tabular}{lll}
\toprule
\textbf{Región} & \textbf{Material efectivo} & \textbf{Resistividad efectiva} \\
\midrule
\texttt{STAR\_VESSEL} & SS316L & $7.5\times10^{-7}~\Omega\,\mathrm{m}$ \\
\texttt{WALL\_INNER} & EUROFER97/RAFM efectivo & $1.0\times10^{-6}~\Omega\,\mathrm{m}$ \\
\texttt{WALL\_OUTER} & EUROFER97/RAFM efectivo & $1.0\times10^{-6}~\Omega\,\mathrm{m}$ \\
Blanket & Región sólida prohibida para plasma & No modelada como conductor continuo en esta etapa \\
\bottomrule
\end{tabular}
\end{table}

Es importante aclarar que la resistividad no altera de forma directa el equilibrio estático en
la corrida actual. La resistividad queda guardada como metadata para diagnósticos, documentación
y posibles extensiones dinámicas donde las corrientes inducidas en pasivos sean relevantes.

\subsection{Corrección del central solenoid y su impacto}
La corrección geométrica más importante fue el reposicionamiento del central solenoid. En la
interpretación previa, el CS estaba demasiado cercano al eje $R=0$. Esto generaba una respuesta
anómala: al aumentar la corriente, el campo tendía a producir compresión radial global, pero no
shaping efectivo de la LCFS.

Al introducir el gap radial observado en el diseño de referencia, aproximadamente
$0.3~\mathrm{m}$, el CS recuperó autoridad sobre la forma del plasma. Esta corrección fue decisiva
para alcanzar simultáneamente $A\approx2$, $\kappa>2$ y $\delta\approx0.47$ sin recurrir a
corrientes imposibles.

El resultado sugiere que el problema principal no era la ausencia de grados de libertad, sino
la posición geométrica de la fuente de campo. Una pequeña diferencia en la posición radial del
winding pack central modificó de forma significativa la sensibilidad del equilibrio.

\subsection{Corrientes del equilibrio baseline}
Las corrientes usadas en el baseline se resumen en la Tabla~\ref{tab:baseline_currents}. Las
corrientes se reportan por familia, comparadas con límites recomendados calculados a partir
de área efectiva y densidad de corriente ingenieril.

\begin{table}[htbp]
\centering
\caption{Corrientes y utilización de bobinas para el equilibrio baseline.}
\label{tab:baseline_currents}
\begin{tabular}{lrrr}
\toprule
\textbf{Familia} & \textbf{$I$ [MA]} & \textbf{$I_{\max}$ [MA]} & \textbf{Utilización [\%]} \\
\midrule
PF1 & $-0.388$ & $15.000$ & $2.59$ \\
PF2 & $-4.374$ & $15.000$ & $29.16$ \\
PF3 & $+1.283$ & $7.500$ & $17.11$ \\
PF4 & $+1.572$ & $5.829$ & $26.97$ \\
PF5 & $+3.753$ & $8.261$ & $45.43$ \\
PF6 & $+6.474$ & $15.302$ & $42.31$ \\
CS total & $-46.655$ & $71.789$ & $64.99$ \\
CS\_MID & $-31.702$ & $32.903$ & $96.35$ \\
CS\_END & $-14.953$ & $38.886$ & $38.45$ \\
\bottomrule
\end{tabular}
\end{table}

La interpretación de CS requiere cuidado. Si \texttt{CS\_MID} y \texttt{CS\_END} se consideran
subregiones de un mismo winding pack, la utilización global del CS es aproximadamente $65\%$.
Sin embargo, la región media local se encuentra cerca de su límite local calculado. En términos
ingenieriles, esto significa que la solución no viola el límite global del CS, pero concentra una
parte importante de la demanda de corriente en la zona media del winding pack.

\subsection{Geometría del equilibrio baseline}
La forma obtenida tiene una frontera LCFS cerrada dentro de la primera pared, con simetría
up-down prácticamente exacta. Sus parámetros geométricos se muestran en la Tabla~\ref{tab:baseline_geometry_detailed}.

\begin{table}[htbp]
\centering
\caption{Parámetros geométricos calculados a partir de la LCFS reconstruida.}
\label{tab:baseline_geometry_detailed}
\begin{tabular}{lc}
\toprule
\textbf{Magnitud} & \textbf{Valor} \\
\midrule
$R_{\mathrm{axis}}$ & $4.397~\mathrm{m}$ \\
$Z_{\mathrm{axis}}$ & $\approx 0~\mathrm{m}$ \\
$R_0$ & $4.208~\mathrm{m}$ \\
$a$ & $2.076~\mathrm{m}$ \\
$A=R_0/a$ & $2.028$ \\
$\kappa$ & $2.163$ \\
$\delta_u$ & $0.469$ \\
$\delta_l$ & $0.469$ \\
Área poloidal & $27.47~\mathrm{m^2}$ \\
Volumen toroidal aproximado & $726.27~\mathrm{m^3}$ \\
$R_{\min},R_{\max}$ & $2.133,6.284~\mathrm{m}$ \\
$Z_{\min},Z_{\max}$ & $-4.489,4.490~\mathrm{m}$ \\
\bottomrule
\end{tabular}
\end{table}

Las definiciones geométricas usadas son las convencionales:

\[
R_0 = \frac{R_{\max}+R_{\min}}{2}, \qquad
 a = \frac{R_{\max}-R_{\min}}{2}, \qquad
 A = \frac{R_0}{a}.
\]

La elongación se calcula como:

\[
\kappa = \frac{Z_{\max}-Z_{\min}}{2a}.
\]

La triangularidad superior e inferior se calculan usando el desplazamiento radial de los puntos
superior e inferior de la LCFS respecto a $R_0$:

\[
\delta_u = \frac{R_0 - R_{Z_{\max}}}{a},\qquad
\delta_l = \frac{R_0 - R_{Z_{\min}}}{a}.
\]

En el baseline, $\delta_u\approx\delta_l$, lo que confirma que la solución es prácticamente
simétrica.

\begin{figure}[htbp]
\centering
\fbox{%
\parbox{0.92\linewidth}{%
\centering
\textbf{[INSERTAR FIGURA]}\\[2mm]
Equilibrio baseline limpio: contornos de $\psi$, LCFS/separatriz, eje magnético, X-points,
\texttt{WALL\_INNER}, \texttt{WALL\_OUTER} y estructuras pasivas con baja opacidad.%
}}
\caption{Equilibrio double-null baseline obtenido con la geometría CAD corregida.}
\label{fig:baseline_equilibrium_clean_placeholder}
\end{figure}

\subsection{Contención y márgenes contra la primera pared}
La LCFS reconstruida se evaluó contra \texttt{WALL\_INNER}, interpretada como first wall/limiter.
El resultado fue:

\begin{table}[htbp]
\centering
\caption{Métricas de contención del equilibrio baseline.}
\label{tab:containment_metrics}
\begin{tabular}{lc}
\toprule
\textbf{Métrica} & \textbf{Valor} \\
\midrule
LCFS dentro de \texttt{WALL\_INNER} & Verdadero \\
Distancia mínima LCFS--\texttt{WALL\_INNER} & $0.0978~\mathrm{m}$ \\
Punto de distancia mínima & $(R,Z)=(4.106,4.085)~\mathrm{m}$ \\
LCFS dentro de \texttt{WALL\_OUTER} & Verdadero \\
Distancia mínima LCFS--\texttt{WALL\_OUTER} & $0.398~\mathrm{m}$ \\
\bottomrule
\end{tabular}
\end{table}

El margen mínimo de casi $10~\mathrm{cm}$ respecto a la primera pared es suficiente para
considerar que el baseline no está simplemente raspando la frontera geométrica. En etapas
posteriores, cuando se incremente presión, esta métrica será una de las primeras a monitorear,
ya que el plasma puede inflarse y reducir rápidamente el margen.

\subsection{Diagnósticos de campo magnético}
Los campos magnéticos se calcularon a partir del flujo poloidal $\psi(R,Z)$ usando:

\[
B_R = -\frac{1}{R}\frac{\partial \psi}{\partial Z},\qquad
B_Z = \frac{1}{R}\frac{\partial \psi}{\partial R},\qquad
B_p = \sqrt{B_R^2+B_Z^2}.
\]

El campo toroidal se aproximó inicialmente como:

\[
B_\phi \approx \frac{f_{\mathrm{vac}}}{R}.
\]

Esta aproximación es adecuada para diagnóstico geométrico inicial, aunque en etapas posteriores
puede reemplazarse por $F(\psi)/R$ si se desea representar explícitamente variaciones internas
del perfil toroidal.

Se generaron dos cortes 1D:

\begin{enumerate}
    \item \textbf{Midplane magnético}: $Z=Z_{\mathrm{axis}}$, barriendo $R$.
    \item \textbf{Eje vertical del plasma}: $R=R_{\mathrm{axis}}$, barriendo $Z$.
\end{enumerate}

En el midplane, $B_R$ funciona como un check de simetría. Para un equilibrio up-down symmetric,
$\partial\psi/\partial Z\approx0$ en $Z=Z_{\mathrm{axis}}$, de modo que $B_R$ debe ser cercano a
cero. En el baseline, $B_R$ en el midplane aparece del orden de $10^{-6}~\mathrm{T}$ o menor,
varias órdenes de magnitud por debajo de $B_Z$, $B_p$ y $B_\phi$. Este residuo se atribuye a
interpolación continua y diferencias finitas, no a un campo radial físicamente relevante.

En el corte vertical, $B_R$ ya no debe ser cero: cambia de signo al cruzar el midplane y muestra
la estructura poloidal asociada a la elongación vertical del plasma. En contraste, $B_\phi$ se
mantiene constante en ese corte porque $R=R_{\mathrm{axis}}$ es constante y se usa
$f_{\mathrm{vac}}/R$.

\begin{figure}[htbp]
\centering
\fbox{%
\parbox{0.92\linewidth}{%
\centering
\textbf{[INSERTAR FIGURA]}\\[2mm]
Perfiles de campo en el midplane: $B_R$, $B_Z$, $B_p$ y $B_\phi$ contra $R$. Deben incluir línea
del eje magnético e intersecciones de LCFS con el midplane.%
}}
\caption{Perfiles magnéticos sobre el midplane magnético.}
\label{fig:midplane_fields_placeholder}
\end{figure}

\begin{figure}[htbp]
\centering
\fbox{%
\parbox{0.92\linewidth}{%
\centering
\textbf{[INSERTAR FIGURA]}\\[2mm]
Perfiles de campo sobre el eje vertical: $B_R$, $B_Z$, $B_p$ y $B_\phi$ contra $Z$ para
$R=R_{\mathrm{axis}}$.%
}}
\caption{Perfiles magnéticos sobre la línea vertical del eje magnético.}
\label{fig:vertical_axis_fields_placeholder}
\end{figure}

\subsection{Perfil de presión y safety factor}
El perfil de presión se obtiene directamente del perfil \texttt{ConstrainPaxisIp}. En el baseline,
la presión cae suavemente desde $p_{\mathrm{axis}}=2.0\times10^3~\mathrm{Pa}$ hasta valores
cercanos a cero en la frontera. Esta presión es baja comparada con metas reactor, pero es
adecuada para validar la rama geométrica.

El perfil de safety factor $q(\psi_N)$ se calculó sobre superficies de flujo normalizadas. El
resultado crece hacia el borde, como se espera en una configuración tokamak. Los valores muy
cercanos a $\psi_N=1$ deben interpretarse con cautela, ya que el cálculo de $q$ cerca de la
separatriz puede ser numéricamente sensible.

\begin{figure}[htbp]
\centering
\fbox{%
\parbox{0.92\linewidth}{%
\centering
\textbf{[INSERTAR FIGURA]}\\[2mm]
Perfil de presión $p(\psi_N)$ del baseline.%
}}
\caption{Perfil de presión normalizado del equilibrio baseline.}
\label{fig:pressure_profile_placeholder}
\end{figure}

\begin{figure}[htbp]
\centering
\fbox{%
\parbox{0.92\linewidth}{%
\centering
\textbf{[INSERTAR FIGURA]}\\[2mm]
Perfil de safety factor $q(\psi_N)$ del baseline.%
}}
\caption{Perfil de safety factor del equilibrio baseline.}
\label{fig:q_profile_placeholder}
\end{figure}

\subsection{Manejo de discrepancias internas de $\psi$}
Durante el diagnóstico apareció el mensaje:

\begin{verbatim}
Discrepancy between psi_func and plasma_psi detected. psi_func has been re-set.
\end{verbatim}

Este mensaje no invalida el equilibrio. Indica que el interpolador interno de $\psi$ no coincidía
exactamente con el arreglo de flujo actualizado, por lo que fue reseteado para recuperar
consistencia. La validación se basa en que, después del reset, los diagnósticos geométricos,
campos, contención y perfiles se calculan de manera coherente. Para verificar estabilidad, se
puede repetir el diagnóstico sin cambiar corrientes; si las métricas principales permanecen
iguales, el warning se considera una corrección interna del objeto de equilibrio.

\subsection{Intentos descartados y su valor metodológico}
Las etapas previas de optimización no se eliminan del historial del proyecto; se documentan
como intentos exploratorios que permitieron identificar errores de modelado. Entre ellas se
incluyen:

\begin{itemize}
    \item optimizaciones directas contra el plasma target CAD original;
    \item refines locales que mejoraban el score pero no la topología;
    \item pruebas con CS exagerado para forzar triangularidad;
    \item segmentaciones adicionales del CS que solo modificaban el equilibrio a corrientes altas;
    \item inverse runs que imponían null points sin resolver la factibilidad geométrica;
    \item continuations en presión que preservaban convergencia numérica pero degradaban la forma.
\end{itemize}

El valor de estas pruebas fue mostrar que el problema no era simplemente de optimización. La
clave estaba en corregir la definición física de la máquina y el posicionamiento del CS. Una vez
hecho esto, apareció una rama de equilibrio mucho más natural.

\subsection{Estado actual del diseño y próximos pasos}
El modelo actual ya está en condiciones de avanzar desde exploración geométrica hacia
continuación física. Antes de aumentar presión, se recomienda congelar este baseline como
referencia y usarlo para comparar futuras ramas. La continuación en presión debe realizarse de
forma gradual:

\[
2\times10^3\ \rightarrow\ 5\times10^3\ \rightarrow\ 10^4\ \rightarrow\ 2\times10^4\ \rightarrow\ 5\times10^4\ \rightarrow\ 10^5~\mathrm{Pa},
\]

monitoreando en cada paso:

\begin{itemize}
    \item contención contra \texttt{WALL\_INNER};
    \item preservación de X-points y topología double-null;
    \item variación de $A$, $\kappa$ y $\delta$;
    \item utilización de corrientes por familia;
    \item evolución de $q(\psi_N)$;
    \item incremento de $\beta$ y energía almacenada cuando esos diagnósticos estén consolidados.
\end{itemize}

El baseline actual no debe sobreoptimizarse manualmente antes de la continuación. Ya cumple la
función principal de esta etapa: demostrar que, con la geometría CAD corregida, el coilset puede
sostener una forma STAR-like double-null confinada y cuantitativamente razonable.

\subsection{Figuras recomendadas para la monografía}
Para la versión final de la monografía se recomienda incluir las siguientes figuras:

\begin{enumerate}
    \item setup CAD actualizado con active coils, STAR\_VESSEL, WALL\_INNER y WALL\_OUTER;
    \item equilibrio baseline limpio;
    \item equilibrio baseline con anotaciones de $R_0$, $a$, X-points y LCFS;
    \item perfiles de campo magnético en midplane;
    \item perfiles de campo magnético sobre eje vertical;
    \item perfil de presión;
    \item perfil de safety factor;
    \item utilización de corrientes de bobinas;
    \item comparación breve entre un intento fallido representativo y el baseline actual.
\end{enumerate}

\begin{figure}[htbp]
\centering
\fbox{%
\parbox{0.92\linewidth}{%
\centering
\textbf{[INSERTAR FIGURA]}\\[2mm]
Resumen visual final sugerido: CAD setup, equilibrio baseline, perfiles magnéticos, $q(\psi_N)$,
presión y utilización de bobinas.%
}}
\caption{Conjunto recomendado de resultados visuales para cerrar la sección del tokamak STAR.}
\label{fig:final_results_panel_placeholder}
\end{figure}

\subsection{Conclusión temporal}
La etapa actual transformó el proyecto de una exploración de formas idealizadas a un pipeline
de diseño ingenieril basado en CAD. Se corrigió la discretización de bobinas, se separaron
estructuras activas y pasivas, se reinterpretó \texttt{WALL\_INNER} como first wall, se eliminó
el blanket pasivo artificial y se corrigió la posición radial del central solenoid. Con estas
modificaciones se obtuvo un equilibrio double-null confinado con $A\approx2.03$, $\kappa\approx2.16$
y $\delta\approx0.47$.

El resultado actual no es todavía el equilibrio final de alta presión, pero sí constituye una
base sólida y reproducible para la siguiente fase: continuation en presión, diagnóstico MHD más
completo y comparación contra objetivos de desempeño reactor.

\end{document}
